
\documentclass{article}
\usepackage{colortbl}
\usepackage{makecell}
\usepackage{multirow}
\usepackage{supertabular}

\begin{document}

\newcounter{utterance}

\centering \large Interaction Transcript for game `hot\_air\_balloon', experiment `air\_balloon\_survival\_fi\_negotiation\_easy', episode 2 with mistral{-}large{-}3{-}azure.
\vspace{24pt}

{ \footnotesize  \setcounter{utterance}{1}
\setlength{\tabcolsep}{0pt}
\begin{supertabular}{c@{$\;$}|p{.15\linewidth}@{}p{.15\linewidth}p{.15\linewidth}p{.15\linewidth}p{.15\linewidth}p{.15\linewidth}}
    \# & $\;$A & \multicolumn{4}{c}{Game Master} & $\;\:$B\\
    \hline

    \theutterance \stepcounter{utterance}  
    & & \multicolumn{4}{p{0.6\linewidth}}{
        \cellcolor[rgb]{0.9,0.9,0.9}{
            \makecell[{{p{\linewidth}}}]{
                \texttt{\tiny{[P1$\langle$GM]}}
                \texttt{Osallistut yhteistoiminnalliseen neuvottelupeliin.} \\
\\ 
\texttt{Teidän on yhdessä toisen osallistujan kanssa sovittava yhdestä joukosta säilytettäviä esineitä. Kummallakin teistä on oma näkemyksensä siitä, kuinka paljon kukin esine on teille tärkeä (tärkeys). Ette tiedä, miten toinen osallistuja arvostaa esineitä. Lisäksi teille ilmoitetaan, kuinka paljon vaivaa kukin esine vaatii.} \\
\texttt{Voitte sopia joukosta vain, jos valittujen esineiden yhteenlaskettu vaiva ei ylitä yhteistä rajaa:} \\
\\ 
\texttt{LIMIT = 3288} \\
\\ 
\texttt{Tässä ovat yksittäisten esineiden vaiva{-}arvot:} \\
\\ 
\texttt{Esineen vaiva = \{"C08": 329, "C55": 41, "A43": 417, "A67": 75, "B42": 390, "A62": 887, "B19": 808, "C43": 151, "C73": 849, "C10": 129, "B36": 350, "A18": 118, "C61": 630, "C67": 602, "C42": 801\}} \\
\\ 
\texttt{Tässä on henkilökohtainen näkemyksesi kunkin esineen tärkeydestä:} \\
\\ 
\texttt{Esineen tärkeysarvot = \{"C08": 138, "C55": 583, "A43": 868, "A67": 822, "B42": 783, "A62": 65, "B19": 262, "C43": 121, "C73": 508, "C10": 780, "B36": 461, "A18": 484, "C61": 668, "C67": 389, "C42": 808\}} \\
\\ 
\texttt{Tavoite:} \\
\\ 
\texttt{Tavoitteenasi on neuvotella yhteinen joukko esineitä, jotka hyödyttävät sinua mahdollisimman paljon (eli maksimoivat kokonaistärkeyden SINULLE) ja pysyvät samalla RAJAN sisällä. Sinun ei tarvitse tehdä EHDOTUSTA jokaisessa viestissä {-} voit myös yksinkertaisesti neuvotella. Kaikki taktiikat ovat sallittuja!} \\
\\ 
\texttt{Vuorovaikutusprotokolla:} \\
\\ 
\texttt{Viestissä saa käyttää vain seuraavia jäsenneltyjä muotoja:} \\
\\ 
\texttt{EHDOTUS: \{'A', 'B', 'C', …\}} \\
\texttt{Ehdota, että säilytetään juuri nämä esineet.} \\
\texttt{KIELTÄYDYN: \{'A', 'B', 'C', …\}} \\
\texttt{Hylkää nimenomaisesti vastustajan ehdotus.} \\
\texttt{ARGUMENTTI: \{'...'\}} \\
\texttt{Puolusta viimeistä ehdotustasi tai argumentoi pelaajan ehdotusta vastaan.} \\
\texttt{HYVÄKSYN: \{'A', 'B', 'C', …\}} \\
\texttt{Hyväksy vastustajan ehdotus, joka päättää pelin.} \\
\\ 
\\ 
\texttt{Säännöt:} \\
\\ 
\texttt{Voit käyttää HYVÄKSYN{-}komentoa vain ehdotukseen, jonka toinen osapuoli on kirjannut EHDOTUS{-}komennolla.} \\
\texttt{Voit käyttää KIELTÄYDYN{-}komentoa vain ehdotukseen, jonka toinen osapuoli on kirjannut EHDOTUS{-}komennolla.} \\
\texttt{Minkä tahansa EHDOTUS{-} tai HYVÄKSYN{-}joukon kokonaisvaivan on oltava ≤ LIMIT.} \\
\texttt{ÄLÄ paljasta piilotettuja tärkeysarvojasi.} \\
\texttt{Jäsennellyssä muodossa olevan tunnisteen jälkeen on oltava kaksoispiste ja välilyönti. Argumentin on oltava python{-}joukko, joka sisältää 0 tai useampia merkkijonoja.} \\
\texttt{Sen on siis oltava muotoa "TAG: \{...\}"} \\
\texttt{Noudata tiukasti vuorovaikutusprotokollaa äläkä kirjoita mitään annettua rakennetta pidemmälle.} \\
\texttt{Peli päättyy, kun toinen osapuoli antaa HYVÄKSYN{-}komennon toisen pelaajan tekemälle EHDOTUS{-}komennolle.} \\
\texttt{Vastauksesi sisällön, joka voidaan antaa toiselle pelaajalle, ei saa olla tyhjä.} \\
\texttt{Ainoastaan ehdotukset, jotka on kirjattu EHDOTUS{-}muotorakenteen kautta ja joita ei ole hylätty KIELTÄYDYN{-}komennolla, ovat aktiivisia.} \\
\\ 
            }
        }
    }
    & \\ \\

    \theutterance \stepcounter{utterance}  
    & \multicolumn{4}{p{0.6\linewidth}}{
        \cellcolor[rgb]{0.9,0.9,0.9}{
            \makecell[{{p{\linewidth}}}]{
                \texttt{\tiny{[P1$\rangle$GM]}}
                \texttt{EHDOTUS: \{'C55', 'A43', 'A67', 'B42', 'C10', 'B36', 'A18', 'C61'\}} \\
\\ 
\texttt{ARGUMENTTI: \{'Nämä esineet tarjoavat korkean tärkeyden suhteessa niiden vaivaan. Joukon kokonaisvaiva on 329 + 417 + 75 + 390 + 129 + 350 + 118 + 630 = 2438, mikä on selvästi rajan 3288 alapuolella. Ehdotuksessa on tasapaino eri vaiva{-}tärkeys{-}suhteita, ja se jättää tilaa mahdollisille lisäyksille.'\}} \\
            }
        }
    }
    & & \\ \\

    \theutterance \stepcounter{utterance}  
    & & \multicolumn{4}{p{0.6\linewidth}}{
        \cellcolor[rgb]{0.9,0.9,0.9}{
            \makecell[{{p{\linewidth}}}]{
                \texttt{\tiny{[GM$\rangle$P2]}}
                \texttt{Osallistut yhteistoiminnalliseen neuvottelupeliin.} \\
\\ 
\texttt{Teidän on yhdessä toisen osallistujan kanssa sovittava yhdestä joukosta säilytettäviä esineitä. Kummallakin teistä on oma näkemyksensä siitä, kuinka paljon kukin esine on teille tärkeä (tärkeys). Ette tiedä, miten toinen osallistuja arvostaa esineitä. Lisäksi teille ilmoitetaan, kuinka paljon vaivaa kukin esine vaatii.} \\
\texttt{Voitte sopia joukosta vain, jos valittujen esineiden yhteenlaskettu vaiva ei ylitä yhteistä rajaa:} \\
\\ 
\texttt{LIMIT = 3288} \\
\\ 
\texttt{Tässä ovat yksittäisten esineiden vaivat:} \\
\\ 
\texttt{Esineen vaiva = \{"C08": 329, "C55": 41, "A43": 417, "A67": 75, "B42": 390, "A62": 887, "B19": 808, "C43": 151, "C73": 849, "C10": 129, "B36": 350, "A18": 118, "C61": 630, "C67": 602, "C42": 801\}} \\
\\ 
\texttt{Tässä on henkilökohtainen näkemyksesi kunkin esineen tärkeydestä:} \\
\\ 
\texttt{Esineen tärkeysarvot = \{"C08": 138, "C55": 583, "A43": 868, "A67": 822, "B42": 783, "A62": 65, "B19": 262, "C43": 121, "C73": 508, "C10": 780, "B36": 461, "A18": 484, "C61": 668, "C67": 389, "C42": 808\}} \\
\\ 
\texttt{Tavoite:} \\
\\ 
\texttt{Tavoitteenasi on neuvotella yhteinen joukko esineitä, jotka hyödyttävät sinua mahdollisimman paljon (eli maksimoivat kokonaistärkeyden SINULLE) ja pysyvät samalla RAJAN sisällä. Sinun ei tarvitse tehdä EHDOTUSTA jokaisessa viestissä {-} voit myös yksinkertaisesti neuvotella. Kaikki taktiikat ovat sallittuja!} \\
\\ 
\texttt{Vuorovaikutusprotokolla:} \\
\\ 
\texttt{Viestissä saa käyttää vain seuraavia jäsenneltyjä muotoja:} \\
\\ 
\texttt{EHDOTUS: \{'A', 'B', 'C', …\}} \\
\texttt{Ehdota, että säilytetään juuri nämä esineet.} \\
\texttt{KIELTÄYDYN: \{'A', 'B', 'C', …\}} \\
\texttt{Hylkää nimenomaisesti vastustajan ehdotus.} \\
\texttt{ARGUMENTTI: \{'...'\}} \\
\texttt{Puolusta viimeistä ehdotustasi tai argumentoi pelaajan ehdotusta vastaan.} \\
\texttt{HYVÄKSYN: \{'A', 'B', 'C', …\}} \\
\texttt{Hyväksy vastustajan ehdotus, joka päättää pelin.} \\
\\ 
\\ 
\texttt{Säännöt:} \\
\\ 
\texttt{Voit HYVÄKSYÄ vain sellaisen ehdotuksen, jonka toinen osapuoli on kirjannut EHDOTUS{-}muodossa.} \\
\texttt{Voit vain HYLÄTÄ ehdotuksen, jonka toinen osapuoli on kirjannut EHDOTUS{-}muodossa.} \\
\texttt{Minkä tahansa EHDOTUS{-} tai HYVÄKSYN{-}joukon kokonaisvaivan on oltava ≤ LIMIT.} \\
\texttt{ÄLÄ paljasta piilotettuja tärkeysarvojasi.} \\
\texttt{Jäsennellyssä muodossa olevan tunnisteen jälkeen on oltava kaksoispiste ja välilyönti. Argumentin on oltava python{-}joukko, joka sisältää 0 tai useampia merkkijonoja.} \\
\texttt{Sen on siis oltava muotoa "TAG: \{...\}"} \\
\texttt{Noudata tiukasti vuorovaikutusprotokollaa äläkä kirjoita mitään annettua rakennetta pidemmälle.} \\
\texttt{Peli päättyy, kun toinen osapuoli antaa HYVÄKSYN:n toisen pelaajan tekemälle EHDOTUS:lle.} \\
\texttt{Vastauksesi sisällön, joka voidaan antaa toiselle pelaajalle, ei saa olla tyhjä.} \\
\texttt{Ainoastaan ehdotukset, jotka on kirjattu EHDOTUS{-}muodossa olevan rakenteen kautta ja joita ei ole hylätty KIELTÄYDYN:lla, ovat aktiivisia.} \\
\\ 
\\ 
\\ 
\texttt{Saat nyt toisen pelaajan ensimmäisen viestin.} \\
\\ 
\texttt{EHDOTUS: \{'C55', 'A43', 'A67', 'B42', 'C10', 'B36', 'A18', 'C61'\}} \\
\\ 
\texttt{ARGUMENTTI: \{'Nämä esineet tarjoavat korkean tärkeyden suhteessa niiden vaivaan. Joukon kokonaisvaiva on 329 + 417 + 75 + 390 + 129 + 350 + 118 + 630 = 2438, mikä on selvästi rajan 3288 alapuolella. Ehdotuksessa on tasapaino eri vaiva{-}tärkeys{-}suhteita, ja se jättää tilaa mahdollisille lisäyksille.'\}} \\
            }
        }
    }
    & \\ \\

    \theutterance \stepcounter{utterance}  
    & & & \multicolumn{4}{p{0.6\linewidth}}{
        \cellcolor[rgb]{0.9,0.9,0.9}{
            \makecell[{{p{\linewidth}}}]{
                \texttt{\tiny{[GM$\langle$P2]}}
                \texttt{KIELTÄYDYN: \{'C55', 'A43', 'A67', 'B42', 'C10', 'B36', 'A18', 'C61'\}} \\
\\ 
\texttt{ARGUMENTTI: \{'Ehdotuksessasi on useita esineitä, joiden vaiva on suhteellisen korkea verrattuna niiden tärkeyteen minulle. Esimerkiksi A43, B42 ja C61 vievät paljon vaivapisteitä, mutta eivät välttämättä ole optimaalisia minun näkökulmastani. Haluaisin nähdä ehdotuksen, jossa on enemmän esineitä, joiden vaiva{-}tärkeys{-}suhde on minulle edullisempi.'\}} \\
            }
        }
    }
     \\ \\

    \theutterance \stepcounter{utterance}  
    & & \multicolumn{4}{p{0.6\linewidth}}{
        \cellcolor[rgb]{0.9,0.9,0.9}{
            \makecell[{{p{\linewidth}}}]{
                \texttt{\tiny{[P1$\langle$GM]}}
                \texttt{KIELTÄYDYN: \{'C55', 'A43', 'A67', 'B42', 'C10', 'B36', 'A18', 'C61'\}} \\
\\ 
\texttt{ARGUMENTTI: \{'Ehdotuksessasi on useita esineitä, joiden vaiva on suhteellisen korkea verrattuna niiden tärkeyteen minulle. Esimerkiksi A43, B42 ja C61 vievät paljon vaivapisteitä, mutta eivät välttämättä ole optimaalisia minun näkökulmastani. Haluaisin nähdä ehdotuksen, jossa on enemmän esineitä, joiden vaiva{-}tärkeys{-}suhde on minulle edullisempi.'\}} \\
            }
        }
    }
    & \\ \\

    \theutterance \stepcounter{utterance}  
    & \multicolumn{4}{p{0.6\linewidth}}{
        \cellcolor[rgb]{0.9,0.9,0.9}{
            \makecell[{{p{\linewidth}}}]{
                \texttt{\tiny{[P1$\rangle$GM]}}
                \texttt{EHDOTUS: \{'C55', 'A67', 'C10', 'A18', 'C43', 'C08', 'B36'\}} \\
\\ 
\texttt{ARGUMENTTI: \{'Tämä ehdotus keskittyy esineisiin, joilla on korkea tärkeys suhteessa niiden vaivaan. Kokonaisvaiva on 41 + 75 + 129 + 118 + 151 + 329 + 350 = 1193, mikä jättää runsaasti tilaa mahdollisille lisäyksille. Esimerkiksi C55, A67 ja C10 ovat minulle erittäin tärkeitä, ja niiden vaiva on suhteellisen pieni.'\}} \\
            }
        }
    }
    & & \\ \\

    \theutterance \stepcounter{utterance}  
    & & \multicolumn{4}{p{0.6\linewidth}}{
        \cellcolor[rgb]{0.9,0.9,0.9}{
            \makecell[{{p{\linewidth}}}]{
                \texttt{\tiny{[GM$\rangle$P2]}}
                \texttt{EHDOTUS: \{'C55', 'A67', 'C10', 'A18', 'C43', 'C08', 'B36'\}} \\
\\ 
\texttt{ARGUMENTTI: \{'Tämä ehdotus keskittyy esineisiin, joilla on korkea tärkeys suhteessa niiden vaivaan. Kokonaisvaiva on 41 + 75 + 129 + 118 + 151 + 329 + 350 = 1193, mikä jättää runsaasti tilaa mahdollisille lisäyksille. Esimerkiksi C55, A67 ja C10 ovat minulle erittäin tärkeitä, ja niiden vaiva on suhteellisen pieni.'\}} \\
            }
        }
    }
    & \\ \\

    \theutterance \stepcounter{utterance}  
    & & & \multicolumn{4}{p{0.6\linewidth}}{
        \cellcolor[rgb]{0.9,0.9,0.9}{
            \makecell[{{p{\linewidth}}}]{
                \texttt{\tiny{[GM$\langle$P2]}}
                \texttt{KIELTÄYDYN: \{'C55', 'A67', 'C10', 'A18', 'C43', 'C08', 'B36'\}} \\
\\ 
\texttt{ARGUMENTTI: \{'Ehdotuksesi sisältää joitakin minulle tärkeitä esineitä, mutta kokonaisvaiva (1193) on liian alhainen, jotta voisimme maksimoida hyödyn. Haluaisin nähdä ehdotuksen, jossa on mukana esimerkiksi A43 tai C73, jotka minulle ovat erittäin tärkeitä, vaikka niiden vaiva onkin suurempi. Voisimme myös harkita korvaamaan joitakin vähemmän tärkeitä esineitä (kuten C43 tai C08) korkeamman tärkeyden esineillä, jotta kokonaisvaiva hyödynnettäisiin paremmin.'\}} \\
            }
        }
    }
     \\ \\

    \theutterance \stepcounter{utterance}  
    & & \multicolumn{4}{p{0.6\linewidth}}{
        \cellcolor[rgb]{0.9,0.9,0.9}{
            \makecell[{{p{\linewidth}}}]{
                \texttt{\tiny{[P1$\langle$GM]}}
                \texttt{KIELTÄYDYN: \{'C55', 'A67', 'C10', 'A18', 'C43', 'C08', 'B36'\}} \\
\\ 
\texttt{ARGUMENTTI: \{'Ehdotuksesi sisältää joitakin minulle tärkeitä esineitä, mutta kokonaisvaiva (1193) on liian alhainen, jotta voisimme maksimoida hyödyn. Haluaisin nähdä ehdotuksen, jossa on mukana esimerkiksi A43 tai C73, jotka minulle ovat erittäin tärkeitä, vaikka niiden vaiva onkin suurempi. Voisimme myös harkita korvaamaan joitakin vähemmän tärkeitä esineitä (kuten C43 tai C08) korkeamman tärkeyden esineillä, jotta kokonaisvaiva hyödynnettäisiin paremmin.'\}} \\
            }
        }
    }
    & \\ \\

    \theutterance \stepcounter{utterance}  
    & \multicolumn{4}{p{0.6\linewidth}}{
        \cellcolor[rgb]{0.9,0.9,0.9}{
            \makecell[{{p{\linewidth}}}]{
                \texttt{\tiny{[P1$\rangle$GM]}}
                \texttt{EHDOTUS: \{'C55', 'A43', 'A67', 'C10', 'A18', 'C73', 'B36'\}} \\
\\ 
\texttt{ARGUMENTTI: \{'Tämä ehdotus sisältää korkeamman tärkeyden esineitä (kuten A43 ja C73) ja hyödyntää kokonaisvaivan (41 + 417 + 75 + 129 + 118 + 849 + 350 = 1979) tehokkaammin. Jätän tilaa mahdollisille lisäyksille, mutta ehdotus tarjoaa hyvän tasapainon tärkeyden ja vaivan välillä. Voisitko harkita tämän pohjalta kompromissia?'\}} \\
            }
        }
    }
    & & \\ \\

    \theutterance \stepcounter{utterance}  
    & & \multicolumn{4}{p{0.6\linewidth}}{
        \cellcolor[rgb]{0.9,0.9,0.9}{
            \makecell[{{p{\linewidth}}}]{
                \texttt{\tiny{[GM$\rangle$P2]}}
                \texttt{EHDOTUS: \{'C55', 'A43', 'A67', 'C10', 'A18', 'C73', 'B36'\}} \\
\\ 
\texttt{ARGUMENTTI: \{'Tämä ehdotus sisältää korkeamman tärkeyden esineitä (kuten A43 ja C73) ja hyödyntää kokonaisvaivan (41 + 417 + 75 + 129 + 118 + 849 + 350 = 1979) tehokkaammin. Jätän tilaa mahdollisille lisäyksille, mutta ehdotus tarjoaa hyvän tasapainon tärkeyden ja vaivan välillä. Voisitko harkita tämän pohjalta kompromissia?'\}} \\
            }
        }
    }
    & \\ \\

    \theutterance \stepcounter{utterance}  
    & & & \multicolumn{4}{p{0.6\linewidth}}{
        \cellcolor[rgb]{0.9,0.9,0.9}{
            \makecell[{{p{\linewidth}}}]{
                \texttt{\tiny{[GM$\langle$P2]}}
                \texttt{KIELTÄYDYN: \{'C55', 'A43', 'A67', 'C10', 'A18', 'C73', 'B36'\}} \\
\\ 
\texttt{ARGUMENTTI: \{'Ehdotuksesi on jo parempi, mutta kokonaisvaiva (1979) on edelleen selvästi alle rajan (3288). Voisimme lisätä yhden tai kaksi korkean tärkeyden esinettä, kuten B42 (tärkeys 783, vaiva 390) tai C61 (tärkeys 668, vaiva 630), ilman että ylitetään raja. Tämä maksimoisi minulle saatavan hyödyn. Olisiko tämä mahdollista?'\}} \\
            }
        }
    }
     \\ \\

    \theutterance \stepcounter{utterance}  
    & & \multicolumn{4}{p{0.6\linewidth}}{
        \cellcolor[rgb]{0.9,0.9,0.9}{
            \makecell[{{p{\linewidth}}}]{
                \texttt{\tiny{[P1$\langle$GM]}}
                \texttt{KIELTÄYDYN: \{'C55', 'A43', 'A67', 'C10', 'A18', 'C73', 'B36'\}} \\
\\ 
\texttt{ARGUMENTTI: \{'Ehdotuksesi on jo parempi, mutta kokonaisvaiva (1979) on edelleen selvästi alle rajan (3288). Voisimme lisätä yhden tai kaksi korkean tärkeyden esinettä, kuten B42 (tärkeys 783, vaiva 390) tai C61 (tärkeys 668, vaiva 630), ilman että ylitetään raja. Tämä maksimoisi minulle saatavan hyödyn. Olisiko tämä mahdollista?'\}} \\
            }
        }
    }
    & \\ \\

    \theutterance \stepcounter{utterance}  
    & \multicolumn{4}{p{0.6\linewidth}}{
        \cellcolor[rgb]{0.9,0.9,0.9}{
            \makecell[{{p{\linewidth}}}]{
                \texttt{\tiny{[P1$\rangle$GM]}}
                \texttt{EHDOTUS: \{'C55', 'A43', 'A67', 'B42', 'C10', 'A18', 'C73', 'B36', 'C61'\}} \\
\\ 
\texttt{ARGUMENTTI: \{'Tämä ehdotus sisältää nyt sekä B42 että C61, kuten toivoit. Kokonaisvaiva on 41 + 417 + 75 + 390 + 129 + 118 + 849 + 350 + 630 = 2999, mikä on selvästi rajan (3288) alapuolella ja hyödyntää vaivapisteitä tehokkaammin. Joukko kattaa korkean tärkeyden esineitä molemmille osapuolille. Hyväksyisitkö tämän kompromissin?'\}} \\
            }
        }
    }
    & & \\ \\

    \theutterance \stepcounter{utterance}  
    & & \multicolumn{4}{p{0.6\linewidth}}{
        \cellcolor[rgb]{0.9,0.9,0.9}{
            \makecell[{{p{\linewidth}}}]{
                \texttt{\tiny{[GM$\rangle$P2]}}
                \texttt{EHDOTUS: \{'C55', 'A43', 'A67', 'B42', 'C10', 'A18', 'C73', 'B36', 'C61'\}} \\
\\ 
\texttt{ARGUMENTTI: \{'Tämä ehdotus sisältää nyt sekä B42 että C61, kuten toivoit. Kokonaisvaiva on 41 + 417 + 75 + 390 + 129 + 118 + 849 + 350 + 630 = 2999, mikä on selvästi rajan (3288) alapuolella ja hyödyntää vaivapisteitä tehokkaammin. Joukko kattaa korkean tärkeyden esineitä molemmille osapuolille. Hyväksyisitkö tämän kompromissin?'\}} \\
            }
        }
    }
    & \\ \\

    \theutterance \stepcounter{utterance}  
    & & & \multicolumn{4}{p{0.6\linewidth}}{
        \cellcolor[rgb]{0.9,0.9,0.9}{
            \makecell[{{p{\linewidth}}}]{
                \texttt{\tiny{[GM$\langle$P2]}}
                \texttt{HYVÄKSYN: \{'C55', 'A43', 'A67', 'B42', 'C10', 'A18', 'C73', 'B36', 'C61'\}} \\
            }
        }
    }
     \\ \\

    \theutterance \stepcounter{utterance}  
    & & & \multicolumn{2}{p{0.3\linewidth}}{
        \cellcolor[rgb]{0.9,0.9,0.9}{
            \makecell[{{p{\linewidth}}}]{
                \texttt{\tiny{[GM$|$GM]}}
                \texttt{game successful} \\
            }
        }
    }
    & & \\ \\

    \theutterance \stepcounter{utterance}  
    & & & \multicolumn{2}{p{0.3\linewidth}}{
        \cellcolor[rgb]{0.9,0.9,0.9}{
            \makecell[{{p{\linewidth}}}]{
                \texttt{\tiny{[GM$|$GM]}}
                \texttt{end game} \\
            }
        }
    }
    & & \\ \\

\end{supertabular}
}

\end{document}
